\documentclass[tikz,border=4pt]{standalone}
% ==========================================================================
% tikz-tcee.tex — Estilos comunes de los gráficos TikZ del temario
% Autor: Víctor Gutiérrez Marcos
%
% Cada gráfico es un fichero standalone en fuentes/<tema>/graficos/:
%
%   \documentclass[tikz,border=4pt]{standalone}
%   % ==========================================================================
% tikz-tcee.tex — Estilos comunes de los gráficos TikZ del temario
% Autor: Víctor Gutiérrez Marcos
%
% Cada gráfico es un fichero standalone en fuentes/<tema>/graficos/:
%
%   \documentclass[tikz,border=4pt]{standalone}
%   % ==========================================================================
% tikz-tcee.tex — Estilos comunes de los gráficos TikZ del temario
% Autor: Víctor Gutiérrez Marcos
%
% Cada gráfico es un fichero standalone en fuentes/<tema>/graficos/:
%
%   \documentclass[tikz,border=4pt]{standalone}
%   \input{../../../latex/tikz-tcee}
%   \begin{document}
%   \begin{tikzpicture}[tcee]
%     \ejes{Q}{P}{6}{5}
%     \draw[curva1] (0.5,4.5) -- (5,0.8) node[etiqueta,right] {$D$};
%     \capa{2}{\draw[curva2] ...;}   % en el vídeo aparece en el paso 2
%   \end{tikzpicture}
%   \end{document}
%
% Paleta: la de la web (morado, dorado, salvia) más tres colores de apoyo
% distinguibles en blanco y negro por el trazo.
% ==========================================================================
\usepackage[T1]{fontenc}
\usepackage{newpxtext,newpxmath}
\usepackage{xcolor}
\usepackage{tikz,pgfplots}
\pgfplotsset{compat=1.18}
\usetikzlibrary{arrows.meta,calc,positioning,patterns,decorations.pathreplacing,intersections,shapes.misc,backgrounds}

\definecolor{tceemorado}{HTML}{5F2987}
\definecolor{tceedorado}{HTML}{B8860B}
\definecolor{tceeverde}{HTML}{3C7D3A}
\definecolor{tceeazul}{HTML}{1F5F99}
\definecolor{tceerojo}{HTML}{B22222}
\definecolor{tceenaranja}{HTML}{E07B28}
\definecolor{tceegris}{HTML}{767171}
\definecolor{tceesalvia}{HTML}{E2EFD9}

% Capas para el vídeo: \capa{n}{…} dibuja su contenido solo a partir del paso n.
% En el PDF, la web y Word se ve todo; generar-video.py compila el gráfico una
% vez por paso con \def\tceepaso{k} para que cada elemento aparezca al narrarlo.
% Vectores con flecha encima, como en el texto (tcee.sty): \vect{x}, \vect{p}
\providecommand{\vect}[1]{\vec{#1}}
\providecommand{\tceepaso}{1000}
\newcommand{\capa}[2]{\ifnum#1>\tceepaso\relax\else#2\fi}
% Gráficos por pasos del mismo tamaño (propuesta del 3.A.29, 11/11/2026): con
% \def\tceefantasma{1} antes del \input (p. ej., x-paso1.tex:
% «\def\tceepaso{3}\def\tceefantasma{1}\input{x.tex}»), las capas posteriores
% al paso se dibujan invisibles y ocupan su sitio. Sin \tceefantasma (el vídeo),
% no cambia nada.
\ifdefined\tceefantasma
  \renewcommand{\capa}[2]{\ifnum#1>\tceepaso\relax\begin{scope}[opacity=0]#2\end{scope}\else#2\fi}
\fi

\tikzset{
  tcee/.style={line cap=round,line join=round,font=\small,>={Stealth[length=2.2mm]}},
  eje/.style={->,thick,black},
  curva1/.style={very thick,tceemorado},
  curva2/.style={very thick,tceeazul},
  curva3/.style={very thick,tceeverde},
  curva4/.style={very thick,tceedorado},
  curvanueva/.style={very thick,tceerojo},
  desplazada/.style={very thick,dashed},
  guia/.style={thin,dashed,tceegris},
  punto/.style={circle,fill=black,inner sep=1.3pt},
  etiqueta/.style={font=\small,inner sep=2pt},
  flecha/.style={->,thick,tceerojo},
  area/.style={fill=tceemorado,fill opacity=0.18,draw=none},
  area2/.style={fill=tceedorado,fill opacity=0.22,draw=none},
}

% \ejes{etiqueta x}{etiqueta y}{largo x}{alto y}
\newcommand{\ejes}[4]{%
  \draw[eje] (0,0) -- (#3,0) node[below left=2pt and 0pt] {#1};
  \draw[eje] (0,0) -- (0,#4) node[left] {#2};
  \node[below left] at (0,0) {$0$};}
% \proyeccion{(x,y)}{etq x}{etq y}: líneas guía a los ejes con etiquetas
\newcommand{\proyeccion}[3]{%
  \draw[guia] let \p1=#1 in (\x1,0) node[below] {#2} |- (0,\y1) node[left] {#3};}

% --------------------------------------------------------------------------
% Paleta «dictamen» (dic…): gráficos de datos de la Guía del dictamen (1D01) y
% de cualquier gráfico de coyuntura que quiera los mismos colores. Un concepto,
% un color en todos los gráficos (CRITERIOS §7, «tres reglas», 24/10/2026);
% tabla completa en fuentes/1D01/_trabajo/paleta-1D01.md. No cambia nada de
% lo anterior: son colores y estilos nuevos.
% Actividad: el PIB (y sus variantes: PIB por persona, PIB mundial), morado.
\definecolor{dicpib}{HTML}{5F2987}
% Demanda interna, gama cálida: consumo (naranjas: privado oscuro, público
% claro), inversión (verde) y existencias (gris con trama, aparte).
\definecolor{dicconspriv}{HTML}{C76A1F}
\definecolor{dicconspub}{HTML}{F5C690}
\definecolor{dicfbcf}{HTML}{2A6B3C}
\definecolor{dicexistencias}{HTML}{B4B4B4}
% Demanda externa y comercio, gama fría (azul claro, que se separa del morado)
\definecolor{dicexterna}{HTML}{6AA6DA}
% Zonas: zona del euro (azul marino), España (rojo), resto de países (grises)
\definecolor{dicze}{HTML}{17406D}
\definecolor{dicespana}{HTML}{C8382E}
\definecolor{dicotro}{HTML}{8A8A8A}
\definecolor{dicotroclaro}{HTML}{B5B5B5}
\definecolor{dicotromasclaro}{HTML}{DADADA}
% Precios: inflación general (negro, continuo) y subyacente (negro,
% discontinuo); componentes: energía (bermellón), alimentos (amarillo) y los
% de la subyacente en azules (bienes industriales claro, servicios oscuro).
\definecolor{dicinflacion}{HTML}{1A1A1A}
\definecolor{dicenergia}{HTML}{D55E00}
\definecolor{dicalimentos}{HTML}{F2D45C}
\definecolor{dicbine}{HTML}{7FB2E0}
\definecolor{dicservicios}{HTML}{1F4E89}
% Trabajo: empleo (verde azulado), horas (el mismo, claro), tasa de actividad
% (verde azulado oscuro), paro (teja); productividad (dorado de la web).
\definecolor{dicempleo}{HTML}{1B7F6B}
\definecolor{dichoras}{HTML}{86CBB8}
\definecolor{dicactividad}{HTML}{0F5247}
\definecolor{dicparo}{HTML}{A23B2A}
\definecolor{dicproductividad}{HTML}{B8860B}
% Hogares: renta disponible (pizarra) y tasa de ahorro (negro); el consumo es
% dicconspriv.
\definecolor{dicrenta}{HTML}{4A5A6A}
\definecolor{dicahorro}{HTML}{1A1A1A}
% Política monetaria: tipos del BCE en azules (depósito, el de la zona del
% euro; operaciones principales, medio; facilidad marginal, claro); balance:
% préstamos a los bancos (dorado claro), resto del activo (gris claro).
\definecolor{dicdeposito}{HTML}{17406D}
\definecolor{dicopf}{HTML}{4F86C6}
\definecolor{dicfmc}{HTML}{9EC5EA}
\definecolor{dicprestamos}{HTML}{D9A520}
\definecolor{dicresto}{HTML}{CFCFCF}
% Ciclo: series adelantadas, coincidentes (el PIB) y retrasadas
\definecolor{dicadelantada}{HTML}{D55E00}
\definecolor{dicretrasada}{HTML}{6AA6DA}
% Referencias: rejilla, previsión, sombreados
\definecolor{dicreferencia}{HTML}{767171}

\tikzset{
  % barra de un componente: \addplot[dicbarra=dicfbcf]
  dicbarra/.style={fill=#1, draw=#1!65!black, line width=0.3pt},
  % existencias: gris con trama, para que se distinga también en gris
  dicbarratrama/.style={fill=#1, draw=#1!65!black, line width=0.3pt,
    postaction={pattern=north east lines, pattern color=white}},
  % línea de una serie; con halo blanco cuando pasa sobre barras
  diclinea/.style={very thick, #1},
  diclineahalo/.style={very thick, #1, preaction={draw=white, line width=3.4pt, solid}},
  % puntos de una serie (capa aparte, después de la línea)
  dicpuntos/.style={only marks, mark=*, mark size=1.8pt,
    mark options={fill=white, draw=#1, line width=0.9pt, solid}},
  % línea del cero o de un umbral, después de las barras
  diccero/.style={thin, black},
  dicumbral/.style={thin, dashed, dicreferencia},
}
% Capas de los gráficos de datos (CRITERIOS §7, «tres reglas»), de abajo
% arriba: cuadrícula → sombreados (dicfondo) → barras y áreas (dicbarras) → ejes, cero y umbrales
% (axis lines) → líneas (main) → puntos (dicpuntos) → rótulos (axis
% descriptions). No basta el orden del código: pgfplots dibuja cuadrícula,
% ejes y marcas en sus propias capas (y area style pone los ejes encima), y
% dentro de un axis los \draw, \fill y \node no admiten capa (van con los
% datos de main); por eso todo lo que se dibuja dentro del axis es un \addplot.
% Uso: \begin{axis}[…, dicapas] (al final de las opciones);
%   barras y áreas: \addplot[on layer=dicbarras, …];
%   sombreados:     \dicsombra[estilo]{x1}{y1}{x2}{y2} (rectángulo, en dicfondo,
%                   bajo las barras y sobre la cuadrícula);
%   cero y umbrales: \dicraya[estilo]{(x1,y1) (x2,y2)} (en axis lines);
%   líneas: \addplot normal (main); puntos: \addplot[on layer=dicpuntos, …]
%   (las marcas de cualquier serie van solas a dicpuntos);
%   rótulos: \dicrotulo[opciones del nodo]{(x,y)}{texto} (en axis descriptions),
%   siempre sin fondo (nada de fill=white: si chocan con una línea, se desplazan;
%   CRITERIOS §7, regla 4).
% Los \capa{n} del vídeo no cambian: siguen envolviendo instrucciones completas.
\pgfdeclarelayer{dicfondo}
\pgfdeclarelayer{dicbarras}
\pgfdeclarelayer{dicpuntos}
\pgfplotsset{
  % (el segundo argumento reparte cuadrícula, ejes y rótulos como el estándar)
  layers/dic/.define layer set={axis background, axis grid, dicfondo, dicbarras,
    axis ticks, axis lines, axis tick labels, pre main, main, dicpuntos,
    axis descriptions, axis foreground}{/pgfplots/layers/standard},
  % leyenda sin caja blanca: toma el fondo del gráfico, también si el gráfico va
  % dentro de un cuadro gris (CRITERIOS §7, regla 4, 26/10/2026)
  dicapas/.style={axis on top=false, set layers=dic, mark layer=dicpuntos,
    legend style={fill=none}},
  dicsuelto/.style={forget plot, stack plots=false, sharp plot, mark=none},
}
\newcommand{\dicraya}[2][]{%
  \addplot[on layer=axis lines, dicsuelto, diccero, #1] coordinates {#2};}
\newcommand{\dicsombra}[5][]{%
  \addplot[on layer=dicfondo, dicsuelto, draw=none, #1]
    coordinates {(#2,#3) (#4,#3) (#4,#5) (#2,#5)};}
\newcommand{\dicrotulo}[3][]{%
  \addplot[on layer=axis descriptions, dicsuelto, draw=none] coordinates {#2}
    node[pos=0, #1] {#3};}

%   \begin{document}
%   \begin{tikzpicture}[tcee]
%     \ejes{Q}{P}{6}{5}
%     \draw[curva1] (0.5,4.5) -- (5,0.8) node[etiqueta,right] {$D$};
%     \capa{2}{\draw[curva2] ...;}   % en el vídeo aparece en el paso 2
%   \end{tikzpicture}
%   \end{document}
%
% Paleta: la de la web (morado, dorado, salvia) más tres colores de apoyo
% distinguibles en blanco y negro por el trazo.
% ==========================================================================
\usepackage[T1]{fontenc}
\usepackage{newpxtext,newpxmath}
\usepackage{xcolor}
\usepackage{tikz,pgfplots}
\pgfplotsset{compat=1.18}
\usetikzlibrary{arrows.meta,calc,positioning,patterns,decorations.pathreplacing,intersections,shapes.misc,backgrounds}

\definecolor{tceemorado}{HTML}{5F2987}
\definecolor{tceedorado}{HTML}{B8860B}
\definecolor{tceeverde}{HTML}{3C7D3A}
\definecolor{tceeazul}{HTML}{1F5F99}
\definecolor{tceerojo}{HTML}{B22222}
\definecolor{tceenaranja}{HTML}{E07B28}
\definecolor{tceegris}{HTML}{767171}
\definecolor{tceesalvia}{HTML}{E2EFD9}

% Capas para el vídeo: \capa{n}{…} dibuja su contenido solo a partir del paso n.
% En el PDF, la web y Word se ve todo; generar-video.py compila el gráfico una
% vez por paso con \def\tceepaso{k} para que cada elemento aparezca al narrarlo.
% Vectores con flecha encima, como en el texto (tcee.sty): \vect{x}, \vect{p}
\providecommand{\vect}[1]{\vec{#1}}
\providecommand{\tceepaso}{1000}
\newcommand{\capa}[2]{\ifnum#1>\tceepaso\relax\else#2\fi}
% Gráficos por pasos del mismo tamaño (propuesta del 3.A.29, 11/11/2026): con
% \def\tceefantasma{1} antes del \input (p. ej., x-paso1.tex:
% «\def\tceepaso{3}\def\tceefantasma{1}\input{x.tex}»), las capas posteriores
% al paso se dibujan invisibles y ocupan su sitio. Sin \tceefantasma (el vídeo),
% no cambia nada.
\ifdefined\tceefantasma
  \renewcommand{\capa}[2]{\ifnum#1>\tceepaso\relax\begin{scope}[opacity=0]#2\end{scope}\else#2\fi}
\fi

\tikzset{
  tcee/.style={line cap=round,line join=round,font=\small,>={Stealth[length=2.2mm]}},
  eje/.style={->,thick,black},
  curva1/.style={very thick,tceemorado},
  curva2/.style={very thick,tceeazul},
  curva3/.style={very thick,tceeverde},
  curva4/.style={very thick,tceedorado},
  curvanueva/.style={very thick,tceerojo},
  desplazada/.style={very thick,dashed},
  guia/.style={thin,dashed,tceegris},
  punto/.style={circle,fill=black,inner sep=1.3pt},
  etiqueta/.style={font=\small,inner sep=2pt},
  flecha/.style={->,thick,tceerojo},
  area/.style={fill=tceemorado,fill opacity=0.18,draw=none},
  area2/.style={fill=tceedorado,fill opacity=0.22,draw=none},
}

% \ejes{etiqueta x}{etiqueta y}{largo x}{alto y}
\newcommand{\ejes}[4]{%
  \draw[eje] (0,0) -- (#3,0) node[below left=2pt and 0pt] {#1};
  \draw[eje] (0,0) -- (0,#4) node[left] {#2};
  \node[below left] at (0,0) {$0$};}
% \proyeccion{(x,y)}{etq x}{etq y}: líneas guía a los ejes con etiquetas
\newcommand{\proyeccion}[3]{%
  \draw[guia] let \p1=#1 in (\x1,0) node[below] {#2} |- (0,\y1) node[left] {#3};}

% --------------------------------------------------------------------------
% Paleta «dictamen» (dic…): gráficos de datos de la Guía del dictamen (1D01) y
% de cualquier gráfico de coyuntura que quiera los mismos colores. Un concepto,
% un color en todos los gráficos (CRITERIOS §7, «tres reglas», 24/10/2026);
% tabla completa en fuentes/1D01/_trabajo/paleta-1D01.md. No cambia nada de
% lo anterior: son colores y estilos nuevos.
% Actividad: el PIB (y sus variantes: PIB por persona, PIB mundial), morado.
\definecolor{dicpib}{HTML}{5F2987}
% Demanda interna, gama cálida: consumo (naranjas: privado oscuro, público
% claro), inversión (verde) y existencias (gris con trama, aparte).
\definecolor{dicconspriv}{HTML}{C76A1F}
\definecolor{dicconspub}{HTML}{F5C690}
\definecolor{dicfbcf}{HTML}{2A6B3C}
\definecolor{dicexistencias}{HTML}{B4B4B4}
% Demanda externa y comercio, gama fría (azul claro, que se separa del morado)
\definecolor{dicexterna}{HTML}{6AA6DA}
% Zonas: zona del euro (azul marino), España (rojo), resto de países (grises)
\definecolor{dicze}{HTML}{17406D}
\definecolor{dicespana}{HTML}{C8382E}
\definecolor{dicotro}{HTML}{8A8A8A}
\definecolor{dicotroclaro}{HTML}{B5B5B5}
\definecolor{dicotromasclaro}{HTML}{DADADA}
% Precios: inflación general (negro, continuo) y subyacente (negro,
% discontinuo); componentes: energía (bermellón), alimentos (amarillo) y los
% de la subyacente en azules (bienes industriales claro, servicios oscuro).
\definecolor{dicinflacion}{HTML}{1A1A1A}
\definecolor{dicenergia}{HTML}{D55E00}
\definecolor{dicalimentos}{HTML}{F2D45C}
\definecolor{dicbine}{HTML}{7FB2E0}
\definecolor{dicservicios}{HTML}{1F4E89}
% Trabajo: empleo (verde azulado), horas (el mismo, claro), tasa de actividad
% (verde azulado oscuro), paro (teja); productividad (dorado de la web).
\definecolor{dicempleo}{HTML}{1B7F6B}
\definecolor{dichoras}{HTML}{86CBB8}
\definecolor{dicactividad}{HTML}{0F5247}
\definecolor{dicparo}{HTML}{A23B2A}
\definecolor{dicproductividad}{HTML}{B8860B}
% Hogares: renta disponible (pizarra) y tasa de ahorro (negro); el consumo es
% dicconspriv.
\definecolor{dicrenta}{HTML}{4A5A6A}
\definecolor{dicahorro}{HTML}{1A1A1A}
% Política monetaria: tipos del BCE en azules (depósito, el de la zona del
% euro; operaciones principales, medio; facilidad marginal, claro); balance:
% préstamos a los bancos (dorado claro), resto del activo (gris claro).
\definecolor{dicdeposito}{HTML}{17406D}
\definecolor{dicopf}{HTML}{4F86C6}
\definecolor{dicfmc}{HTML}{9EC5EA}
\definecolor{dicprestamos}{HTML}{D9A520}
\definecolor{dicresto}{HTML}{CFCFCF}
% Ciclo: series adelantadas, coincidentes (el PIB) y retrasadas
\definecolor{dicadelantada}{HTML}{D55E00}
\definecolor{dicretrasada}{HTML}{6AA6DA}
% Referencias: rejilla, previsión, sombreados
\definecolor{dicreferencia}{HTML}{767171}

\tikzset{
  % barra de un componente: \addplot[dicbarra=dicfbcf]
  dicbarra/.style={fill=#1, draw=#1!65!black, line width=0.3pt},
  % existencias: gris con trama, para que se distinga también en gris
  dicbarratrama/.style={fill=#1, draw=#1!65!black, line width=0.3pt,
    postaction={pattern=north east lines, pattern color=white}},
  % línea de una serie; con halo blanco cuando pasa sobre barras
  diclinea/.style={very thick, #1},
  diclineahalo/.style={very thick, #1, preaction={draw=white, line width=3.4pt, solid}},
  % puntos de una serie (capa aparte, después de la línea)
  dicpuntos/.style={only marks, mark=*, mark size=1.8pt,
    mark options={fill=white, draw=#1, line width=0.9pt, solid}},
  % línea del cero o de un umbral, después de las barras
  diccero/.style={thin, black},
  dicumbral/.style={thin, dashed, dicreferencia},
}
% Capas de los gráficos de datos (CRITERIOS §7, «tres reglas»), de abajo
% arriba: cuadrícula → sombreados (dicfondo) → barras y áreas (dicbarras) → ejes, cero y umbrales
% (axis lines) → líneas (main) → puntos (dicpuntos) → rótulos (axis
% descriptions). No basta el orden del código: pgfplots dibuja cuadrícula,
% ejes y marcas en sus propias capas (y area style pone los ejes encima), y
% dentro de un axis los \draw, \fill y \node no admiten capa (van con los
% datos de main); por eso todo lo que se dibuja dentro del axis es un \addplot.
% Uso: \begin{axis}[…, dicapas] (al final de las opciones);
%   barras y áreas: \addplot[on layer=dicbarras, …];
%   sombreados:     \dicsombra[estilo]{x1}{y1}{x2}{y2} (rectángulo, en dicfondo,
%                   bajo las barras y sobre la cuadrícula);
%   cero y umbrales: \dicraya[estilo]{(x1,y1) (x2,y2)} (en axis lines);
%   líneas: \addplot normal (main); puntos: \addplot[on layer=dicpuntos, …]
%   (las marcas de cualquier serie van solas a dicpuntos);
%   rótulos: \dicrotulo[opciones del nodo]{(x,y)}{texto} (en axis descriptions),
%   siempre sin fondo (nada de fill=white: si chocan con una línea, se desplazan;
%   CRITERIOS §7, regla 4).
% Los \capa{n} del vídeo no cambian: siguen envolviendo instrucciones completas.
\pgfdeclarelayer{dicfondo}
\pgfdeclarelayer{dicbarras}
\pgfdeclarelayer{dicpuntos}
\pgfplotsset{
  % (el segundo argumento reparte cuadrícula, ejes y rótulos como el estándar)
  layers/dic/.define layer set={axis background, axis grid, dicfondo, dicbarras,
    axis ticks, axis lines, axis tick labels, pre main, main, dicpuntos,
    axis descriptions, axis foreground}{/pgfplots/layers/standard},
  % leyenda sin caja blanca: toma el fondo del gráfico, también si el gráfico va
  % dentro de un cuadro gris (CRITERIOS §7, regla 4, 26/10/2026)
  dicapas/.style={axis on top=false, set layers=dic, mark layer=dicpuntos,
    legend style={fill=none}},
  dicsuelto/.style={forget plot, stack plots=false, sharp plot, mark=none},
}
\newcommand{\dicraya}[2][]{%
  \addplot[on layer=axis lines, dicsuelto, diccero, #1] coordinates {#2};}
\newcommand{\dicsombra}[5][]{%
  \addplot[on layer=dicfondo, dicsuelto, draw=none, #1]
    coordinates {(#2,#3) (#4,#3) (#4,#5) (#2,#5)};}
\newcommand{\dicrotulo}[3][]{%
  \addplot[on layer=axis descriptions, dicsuelto, draw=none] coordinates {#2}
    node[pos=0, #1] {#3};}

%   \begin{document}
%   \begin{tikzpicture}[tcee]
%     \ejes{Q}{P}{6}{5}
%     \draw[curva1] (0.5,4.5) -- (5,0.8) node[etiqueta,right] {$D$};
%     \capa{2}{\draw[curva2] ...;}   % en el vídeo aparece en el paso 2
%   \end{tikzpicture}
%   \end{document}
%
% Paleta: la de la web (morado, dorado, salvia) más tres colores de apoyo
% distinguibles en blanco y negro por el trazo.
% ==========================================================================
\usepackage[T1]{fontenc}
\usepackage{newpxtext,newpxmath}
\usepackage{xcolor}
\usepackage{tikz,pgfplots}
\pgfplotsset{compat=1.18}
\usetikzlibrary{arrows.meta,calc,positioning,patterns,decorations.pathreplacing,intersections,shapes.misc,backgrounds}

\definecolor{tceemorado}{HTML}{5F2987}
\definecolor{tceedorado}{HTML}{B8860B}
\definecolor{tceeverde}{HTML}{3C7D3A}
\definecolor{tceeazul}{HTML}{1F5F99}
\definecolor{tceerojo}{HTML}{B22222}
\definecolor{tceenaranja}{HTML}{E07B28}
\definecolor{tceegris}{HTML}{767171}
\definecolor{tceesalvia}{HTML}{E2EFD9}

% Capas para el vídeo: \capa{n}{…} dibuja su contenido solo a partir del paso n.
% En el PDF, la web y Word se ve todo; generar-video.py compila el gráfico una
% vez por paso con \def\tceepaso{k} para que cada elemento aparezca al narrarlo.
% Vectores con flecha encima, como en el texto (tcee.sty): \vect{x}, \vect{p}
\providecommand{\vect}[1]{\vec{#1}}
\providecommand{\tceepaso}{1000}
\newcommand{\capa}[2]{\ifnum#1>\tceepaso\relax\else#2\fi}
% Gráficos por pasos del mismo tamaño (propuesta del 3.A.29, 11/11/2026): con
% \def\tceefantasma{1} antes del \input (p. ej., x-paso1.tex:
% «\def\tceepaso{3}\def\tceefantasma{1}\input{x.tex}»), las capas posteriores
% al paso se dibujan invisibles y ocupan su sitio. Sin \tceefantasma (el vídeo),
% no cambia nada.
\ifdefined\tceefantasma
  \renewcommand{\capa}[2]{\ifnum#1>\tceepaso\relax\begin{scope}[opacity=0]#2\end{scope}\else#2\fi}
\fi

\tikzset{
  tcee/.style={line cap=round,line join=round,font=\small,>={Stealth[length=2.2mm]}},
  eje/.style={->,thick,black},
  curva1/.style={very thick,tceemorado},
  curva2/.style={very thick,tceeazul},
  curva3/.style={very thick,tceeverde},
  curva4/.style={very thick,tceedorado},
  curvanueva/.style={very thick,tceerojo},
  desplazada/.style={very thick,dashed},
  guia/.style={thin,dashed,tceegris},
  punto/.style={circle,fill=black,inner sep=1.3pt},
  etiqueta/.style={font=\small,inner sep=2pt},
  flecha/.style={->,thick,tceerojo},
  area/.style={fill=tceemorado,fill opacity=0.18,draw=none},
  area2/.style={fill=tceedorado,fill opacity=0.22,draw=none},
}

% \ejes{etiqueta x}{etiqueta y}{largo x}{alto y}
\newcommand{\ejes}[4]{%
  \draw[eje] (0,0) -- (#3,0) node[below left=2pt and 0pt] {#1};
  \draw[eje] (0,0) -- (0,#4) node[left] {#2};
  \node[below left] at (0,0) {$0$};}
% \proyeccion{(x,y)}{etq x}{etq y}: líneas guía a los ejes con etiquetas
\newcommand{\proyeccion}[3]{%
  \draw[guia] let \p1=#1 in (\x1,0) node[below] {#2} |- (0,\y1) node[left] {#3};}

% --------------------------------------------------------------------------
% Paleta «dictamen» (dic…): gráficos de datos de la Guía del dictamen (1D01) y
% de cualquier gráfico de coyuntura que quiera los mismos colores. Un concepto,
% un color en todos los gráficos (CRITERIOS §7, «tres reglas», 24/10/2026);
% tabla completa en fuentes/1D01/_trabajo/paleta-1D01.md. No cambia nada de
% lo anterior: son colores y estilos nuevos.
% Actividad: el PIB (y sus variantes: PIB por persona, PIB mundial), morado.
\definecolor{dicpib}{HTML}{5F2987}
% Demanda interna, gama cálida: consumo (naranjas: privado oscuro, público
% claro), inversión (verde) y existencias (gris con trama, aparte).
\definecolor{dicconspriv}{HTML}{C76A1F}
\definecolor{dicconspub}{HTML}{F5C690}
\definecolor{dicfbcf}{HTML}{2A6B3C}
\definecolor{dicexistencias}{HTML}{B4B4B4}
% Demanda externa y comercio, gama fría (azul claro, que se separa del morado)
\definecolor{dicexterna}{HTML}{6AA6DA}
% Zonas: zona del euro (azul marino), España (rojo), resto de países (grises)
\definecolor{dicze}{HTML}{17406D}
\definecolor{dicespana}{HTML}{C8382E}
\definecolor{dicotro}{HTML}{8A8A8A}
\definecolor{dicotroclaro}{HTML}{B5B5B5}
\definecolor{dicotromasclaro}{HTML}{DADADA}
% Precios: inflación general (negro, continuo) y subyacente (negro,
% discontinuo); componentes: energía (bermellón), alimentos (amarillo) y los
% de la subyacente en azules (bienes industriales claro, servicios oscuro).
\definecolor{dicinflacion}{HTML}{1A1A1A}
\definecolor{dicenergia}{HTML}{D55E00}
\definecolor{dicalimentos}{HTML}{F2D45C}
\definecolor{dicbine}{HTML}{7FB2E0}
\definecolor{dicservicios}{HTML}{1F4E89}
% Trabajo: empleo (verde azulado), horas (el mismo, claro), tasa de actividad
% (verde azulado oscuro), paro (teja); productividad (dorado de la web).
\definecolor{dicempleo}{HTML}{1B7F6B}
\definecolor{dichoras}{HTML}{86CBB8}
\definecolor{dicactividad}{HTML}{0F5247}
\definecolor{dicparo}{HTML}{A23B2A}
\definecolor{dicproductividad}{HTML}{B8860B}
% Hogares: renta disponible (pizarra) y tasa de ahorro (negro); el consumo es
% dicconspriv.
\definecolor{dicrenta}{HTML}{4A5A6A}
\definecolor{dicahorro}{HTML}{1A1A1A}
% Política monetaria: tipos del BCE en azules (depósito, el de la zona del
% euro; operaciones principales, medio; facilidad marginal, claro); balance:
% préstamos a los bancos (dorado claro), resto del activo (gris claro).
\definecolor{dicdeposito}{HTML}{17406D}
\definecolor{dicopf}{HTML}{4F86C6}
\definecolor{dicfmc}{HTML}{9EC5EA}
\definecolor{dicprestamos}{HTML}{D9A520}
\definecolor{dicresto}{HTML}{CFCFCF}
% Ciclo: series adelantadas, coincidentes (el PIB) y retrasadas
\definecolor{dicadelantada}{HTML}{D55E00}
\definecolor{dicretrasada}{HTML}{6AA6DA}
% Referencias: rejilla, previsión, sombreados
\definecolor{dicreferencia}{HTML}{767171}

\tikzset{
  % barra de un componente: \addplot[dicbarra=dicfbcf]
  dicbarra/.style={fill=#1, draw=#1!65!black, line width=0.3pt},
  % existencias: gris con trama, para que se distinga también en gris
  dicbarratrama/.style={fill=#1, draw=#1!65!black, line width=0.3pt,
    postaction={pattern=north east lines, pattern color=white}},
  % línea de una serie; con halo blanco cuando pasa sobre barras
  diclinea/.style={very thick, #1},
  diclineahalo/.style={very thick, #1, preaction={draw=white, line width=3.4pt, solid}},
  % puntos de una serie (capa aparte, después de la línea)
  dicpuntos/.style={only marks, mark=*, mark size=1.8pt,
    mark options={fill=white, draw=#1, line width=0.9pt, solid}},
  % línea del cero o de un umbral, después de las barras
  diccero/.style={thin, black},
  dicumbral/.style={thin, dashed, dicreferencia},
}
% Capas de los gráficos de datos (CRITERIOS §7, «tres reglas»), de abajo
% arriba: cuadrícula → sombreados (dicfondo) → barras y áreas (dicbarras) → ejes, cero y umbrales
% (axis lines) → líneas (main) → puntos (dicpuntos) → rótulos (axis
% descriptions). No basta el orden del código: pgfplots dibuja cuadrícula,
% ejes y marcas en sus propias capas (y area style pone los ejes encima), y
% dentro de un axis los \draw, \fill y \node no admiten capa (van con los
% datos de main); por eso todo lo que se dibuja dentro del axis es un \addplot.
% Uso: \begin{axis}[…, dicapas] (al final de las opciones);
%   barras y áreas: \addplot[on layer=dicbarras, …];
%   sombreados:     \dicsombra[estilo]{x1}{y1}{x2}{y2} (rectángulo, en dicfondo,
%                   bajo las barras y sobre la cuadrícula);
%   cero y umbrales: \dicraya[estilo]{(x1,y1) (x2,y2)} (en axis lines);
%   líneas: \addplot normal (main); puntos: \addplot[on layer=dicpuntos, …]
%   (las marcas de cualquier serie van solas a dicpuntos);
%   rótulos: \dicrotulo[opciones del nodo]{(x,y)}{texto} (en axis descriptions),
%   siempre sin fondo (nada de fill=white: si chocan con una línea, se desplazan;
%   CRITERIOS §7, regla 4).
% Los \capa{n} del vídeo no cambian: siguen envolviendo instrucciones completas.
\pgfdeclarelayer{dicfondo}
\pgfdeclarelayer{dicbarras}
\pgfdeclarelayer{dicpuntos}
\pgfplotsset{
  % (el segundo argumento reparte cuadrícula, ejes y rótulos como el estándar)
  layers/dic/.define layer set={axis background, axis grid, dicfondo, dicbarras,
    axis ticks, axis lines, axis tick labels, pre main, main, dicpuntos,
    axis descriptions, axis foreground}{/pgfplots/layers/standard},
  % leyenda sin caja blanca: toma el fondo del gráfico, también si el gráfico va
  % dentro de un cuadro gris (CRITERIOS §7, regla 4, 26/10/2026)
  dicapas/.style={axis on top=false, set layers=dic, mark layer=dicpuntos,
    legend style={fill=none}},
  dicsuelto/.style={forget plot, stack plots=false, sharp plot, mark=none},
}
\newcommand{\dicraya}[2][]{%
  \addplot[on layer=axis lines, dicsuelto, diccero, #1] coordinates {#2};}
\newcommand{\dicsombra}[5][]{%
  \addplot[on layer=dicfondo, dicsuelto, draw=none, #1]
    coordinates {(#2,#3) (#4,#3) (#4,#5) (#2,#5)};}
\newcommand{\dicrotulo}[3][]{%
  \addplot[on layer=axis descriptions, dicsuelto, draw=none] coordinates {#2}
    node[pos=0, #1] {#3};}

% Colores de las cajas del tema (los de tcee-comun.sty)
\definecolor{tceemoradopalido}{HTML}{F3EEF7}
\definecolor{tceeazulpalido}{HTML}{DDEBF7}
\begin{document}
\hyphenpenalty=10000 \exhyphenpenalty=10000
% La sucesión de cuentas del SEC 2010 y sus saldos (elaboración propia a
% partir del SEC 2010, cap. 8). Cascada vertical: cada flecha lleva el saldo
% que pasa a la cuenta siguiente (negrita y morado) y a la derecha de cada caja
% va el agregado del total de la economía. Colores por bloque: cuentas
% corrientes en morado, acumulación en salvia y verde, balances en azul; la
% cuenta de bienes y servicios y el resto del mundo, en gris.
% Es el gráfico de referencia de la ficha identidades-contabilidad-nacional.
% Capas: cuenta 0 -> I -> II.1.1 -> II.1.2 -> II.2 -> II.3 -> II.4 -> llave
% corrientes -> III.1 -> III.2 -> III.3 -> llave acumulación -> IV -> llave
% balances -> resto del mundo.
\begin{tikzpicture}[tcee, font=\footnotesize,
  cuenta/.style={draw=#1, line width=0.8pt, rounded corners=2pt, align=center,
    text width=4.7cm, minimum height=0.8cm, inner sep=3pt},
  corriente/.style={cuenta=tceemorado, fill=tceemoradopalido},
  acumulacion/.style={cuenta=tceeverde, fill=tceesalvia},
  balance/.style={cuenta=tceeazul, fill=tceeazulpalido},
  cuentabys/.style={cuenta=tceegris, fill=tceegris!12},
  paso/.style={->, thick, tceegris!80!black},
  saldo/.style={font=\footnotesize\bfseries, text=tceemorado, anchor=west, align=left,
    inner sep=2pt},
  agregado/.style={font=\footnotesize, anchor=west, align=left, inner sep=2pt},
  total/.style={font=\scriptsize, text=tceegris!80!black},
  llave/.style={decorate, decoration={brace, amplitude=6pt, mirror}, thick, tceegris!80!black},
  bloque/.style={rotate=90, font=\footnotesize\itshape, text=tceegris!80!black},
  rdm/.style={draw=tceegris, line width=0.6pt, fill=tceegris!12, rounded corners=2pt,
    align=center, text width=3.1cm, inner sep=3pt, minimum height=0.8cm}]
  % Altura de cada cuenta (centro de la caja); las cajas miden 0,8 de alto
  \foreach \n/\y in {c0/0, c1/-1.6, c211/-3.2, c212/-4.8, c22/-6.4, c23/-8.0,
      c24/-9.6, c31/-11.2, c32/-13.0, c33/-14.8, c4/-16.6}
    \coordinate (\n) at (0,\y);
  % Marco fijo: todas las capas del vídeo con el mismo tamaño
  \path (-3.9,-18.6) rectangle (11.6,2.1);
  % 1: cuenta de bienes y servicios, la identidad de partida
  \capa{1}{\node[cuentabys] at (c0) {\textbf{0.} Cuenta de bienes y servicios};}
  \capa{1}{\node[agregado, text=tceegris!80!black] at ($(c0)+(2.75,0)$) {recursos $=$ empleos (sin saldo)};}
  % 2: producción; saldo VAB; agregado PIB
  \capa{2}{\draw[paso] ($(c0)+(0,-0.4)$) -- ($(c1)+(0,0.4)$);}
  \capa{2}{\node[corriente] at (c1) {\textbf{I.} Producción};}
  \capa{2}{\draw[paso] ($(c1)+(0,-0.4)$) -- node[saldo, right] {VAB} ($(c211)+(0,0.4)$);}
  \capa{2}{\node[agregado] at ($(c1)+(2.75,0)$) {\textcolor{tceegris!80!black}{total:} PIB};}
  % 3: explotación; saldo EBE y renta mixta
  \capa{3}{\node[corriente] at (c211) {\textbf{II.1.1} Explotación};}
  \capa{3}{\draw[paso] ($(c211)+(0,-0.4)$) -- node[saldo, right] {EBE y renta mixta} ($(c212)+(0,0.4)$);}
  % 4: asignación de la renta primaria; saldo de rentas primarias; RNB
  \capa{4}{\node[corriente] at (c212) {\textbf{II.1.2} Asignación de la renta primaria};}
  \capa{4}{\draw[paso] ($(c212)+(0,-0.4)$) -- node[saldo, right] {saldo de rentas primarias} ($(c22)+(0,0.4)$);}
  \capa{4}{\node[agregado] at ($(c212)+(2.75,0)$) {\textcolor{tceegris!80!black}{total:} RNB};}
  % 5: distribución secundaria; renta disponible; RNDB
  \capa{5}{\node[corriente] at (c22) {\textbf{II.2} Distribución secundaria de la renta};}
  \capa{5}{\draw[paso] ($(c22)+(0,-0.4)$) -- node[saldo, right] {renta disponible} ($(c23)+(0,0.4)$);}
  \capa{5}{\node[agregado] at ($(c22)+(2.75,0)$) {\textcolor{tceegris!80!black}{total:} RNDB};}
  % 6: redistribución en especie; renta disponible ajustada
  \capa{6}{\node[corriente] at (c23) {\textbf{II.3} Redistribución de la renta en especie};}
  \capa{6}{\draw[paso] ($(c23)+(0,-0.4)$) -- node[saldo, right] {renta disponible ajustada} ($(c24)+(0,0.4)$);}
  \capa{6}{\node[agregado] at ($(c23)+(2.75,0)$) {\textcolor{tceegris!80!black}{total:} $=$ RNDB};}
  % 7: utilización de la renta; ahorro
  \capa{7}{\node[corriente] at (c24) {\textbf{II.4} Utilización de la renta};}
  \capa{7}{\draw[paso] ($(c24)+(0,-0.4)$) -- node[saldo, right] {ahorro} ($(c31)+(0,0.4)$);}
  \capa{7}{\node[agregado] at ($(c24)+(2.75,0)$) {\textcolor{tceegris!80!black}{total:} $S=\mathrm{RNDB}-C-G$};}
  % 8: llave de las cuentas corrientes
  \capa{8}{\draw[llave] ($(c1)+(-2.75,0.4)$) -- ($(c24)+(-2.75,-0.4)$)
    node[midway, left=8pt, bloque, anchor=south] {cuentas corrientes};}
  % 9: capital; capacidad o necesidad de financiación
  \capa{9}{\node[acumulacion] at (c31) {\textbf{III.1} Capital};}
  \capa{9}{\draw[paso] ($(c31)+(0,-0.4)$) -- node[saldo, right]
    {capacidad ($+$) o necesidad ($-$)\\ de financiación (B.9)} ($(c32)+(0,0.4)$);}
  % 10: financiera; B.9 = VNA - VNP
  \capa{10}{\node[acumulacion] at (c32) {\textbf{III.2} Financiera};}
  \capa{10}{\draw[paso] ($(c32)+(0,-0.4)$) -- node[saldo, right]
    {$\mathbf{B.9=VNA-VNP}$\\ \mdseries\textcolor{tceegris!80!black}{(salvo la discrepancia estadística)}} ($(c33)+(0,0.4)$);}
  % 11: otras variaciones de los activos
  \capa{11}{\node[acumulacion] at (c33) {\textbf{III.3} Otras variaciones de los activos};}
  \capa{11}{\draw[paso] ($(c33)+(0,-0.4)$) -- node[saldo, right]
    {revalorizaciones y otros\\ cambios de volumen} ($(c4)+(0,0.4)$);}
  % 12: llave de las cuentas de acumulación
  \capa{12}{\draw[llave] ($(c31)+(-2.75,0.4)$) -- ($(c33)+(-2.75,-0.4)$)
    node[midway, left=8pt, bloque, anchor=south] {cuentas de acumulación};}
  % 13: balances; patrimonio neto y vuelta al balance inicial
  \capa{13}{\node[balance] at (c4) {\textbf{IV.} Balances};}
  \capa{13}{\draw[paso] ($(c4)+(0,-0.4)$) -- node[saldo, right] {patrimonio neto} ($(c4)+(0,-1.2)$);}
  \capa{13}{\draw[paso, tceeazul] ($(c4)+(0,-1.2)$) -- ++(0,-0.3) -| ($(c4)+(2.75,0)$) -- ($(c4)+(2.42,0)$);}
  \capa{13}{\node[agregado, text=tceeazul] at ($(c4)+(2.9,-0.75)$)
    {balance final $=$\\ balance inicial $+$ variaciones};}
  % 14: llave de los balances
  \capa{14}{\draw[llave] ($(c4)+(-2.75,0.4)$) -- ($(c4)+(-2.75,-0.4)$)
    node[midway, left=8pt, bloque, anchor=south] {balances};}
  % 15: el resto del mundo, con el signo cambiado
  \capa{15}{\draw[thin, dashed, tceegris] (7.9,2.0) -- (7.9,-13.6);}
  \capa{15}{\node[font=\small\bfseries, align=center] at (9.85,1.5)
    {Resto del mundo (S.2)};}
  \capa{15}{\node[font=\scriptsize, text=tceegris!80!black] at (9.85,0.98)
    {{\normalsize$\boldsymbol{-}$}\, con el signo cambiado};}
  \capa{15}{\node[rdm] at (9.85,0) {\textbf{V.I} bienes y servicios:\\ $M-X$};}
  \capa{15}{\node[rdm, minimum height=2.4cm] at (9.85,-5.6) {\textbf{V.II} rentas primarias y secundarias:\\ $-\mathrm{CC}$};}
  \capa{15}{\node[rdm] at (9.85,-11.2) {\textbf{V.III.1} capital:\\ $-(\mathrm{CC}+\mathrm{CK})$};}
  \capa{15}{\node[rdm] at (9.85,-13.0) {\textbf{V.III.2} financiera};}
\end{tikzpicture}
\end{document}
